% archon:covers MR2223407ConstructionHilbertQuot/SemicontinuityBaseChange.lean
% Nitsure §3 (Semi-Continuity and Base-Change). Source file:
% references/MR2223407-construction-hilbert-quot.tex.
\chapter{Semicontinuity and Base Change}
\label{chap:semicontinuity}

This chapter develops the cohomological base-change machinery of Nitsure §3. It
provides two things used pervasively in the construction of the Quot scheme:
(a) a recognition criterion for flatness of a family from local freeness of the
direct images $\pi_*\Ff(r)$, and (b) the way cohomology of a flat coherent sheaf
varies in families. The deepest input is the \emph{Grothendieck complex}
(EGA III §7.7), recorded in Chapter~\ref{chap:basic} as the anchor
\cref{thm:cohomology-base-change}; here it is restated in the project's notation
and its consequences (semicontinuity, base change for flat sheaves) are spelled
out.

\section{Base change without flatness}

The following lemma on base change does not need any flatness hypothesis. The
price paid is that the integer $r_0$ may depend on $\phi$.

\begin{lemma}\leanok
[Base change without flatness]
  \label{lem:base-change-without-flatness}
  \lean{MR2223407ConstructionHilbertQuot.base_change_without_flatness}
  \uses{def:projective-space, thm:serre-vanishing, lem:five-lemma}
  % SOURCE: Nitsure, §3, lines 1278--1288 (read from references/MR2223407-construction-hilbert-quot.tex)
  % SOURCE QUOTE: "Let $\phi : T\to S$ be a morphism of noetherian schemes,
  % let $\F$ a coherent sheaf on $\PP^n_S$, and
  % let $\F_T$ denote the pull-back of
  % $\F$ under the induced morphism $\PP^n_T\to \PP^n_S$. Let
  % $\pi_S : \PP^n_S \to S$ and $\pi_T : \PP^n_T \to T$
  % denote the projections. Then there exists an integer
  % $r_0$ such that the base-change homomorphism
  % $$\phi^* {\pi_S}_*\, \F(r) \to {\pi_T}_* \, \F_T(r)$$
  % is an isomorphism for all $r\ge r_0$."
  \textit{Source: Nitsure §3 (Base-change without Flatness).}
  Let $\phi\colon T\to S$ be a morphism of noetherian schemes, let $\Ff$ be a
  coherent sheaf on $\PP^n_S$, and let $\Ff_T$ denote the pull-back of $\Ff$
  under the induced morphism $\PP^n_T\to\PP^n_S$. Let
  $\pi_S\colon \PP^n_S\to S$ and $\pi_T\colon \PP^n_T\to T$ denote the
  projections. Then there exists an integer $r_0$ such that the base-change
  homomorphism
  \[
    \phi^*{\pi_S}_*\,\Ff(r)\;\longrightarrow\;{\pi_T}_*\,\Ff_T(r)
  \]
  is an isomorphism for all $r\ge r_0$.
\end{lemma}

% SOURCE QUOTE PROOF: "As base-change holds for open embeddings, using a finite affine open
% cover $U_i$ of $S$ and a finite affine open cover $V_{i,j}$ of each
% $\phi^{-1}(U_i)$ (which is possible by noetherian hypothesis),
% it is enough to consider the case where $S$ and
% $T$ are affine.
% Note that for all integers $i$, the base-change homomorphism
% $$\phi^*{\pi_S}_*{\mathcal O}_{\PP^n_S}(i) \to {\pi_T}_* {\mathcal O}_{\PP^n_T}(i)$$
% is an isomorphism. [...]
% As $S$ is noetherian and affine, there exists an exact sequence
% $$\oplus^p {\mathcal O}_{\PP^n_S}(a)\, \stackrel{u}{\to}  \,
% \oplus^q {\mathcal O}_{\PP^n_S}(b) \,\stackrel{v}{\to} \,\F \,\to \,0$$
% for some integers $a$, $b$, $p\ge 0$, $q\ge 0$. [...]
% Let $\G = \ker(v)$ and let $\H = \ker(v_T)$.
% For any integer $r$, we get exact sequences
% $${\pi_S}_* \oplus^p {\mathcal O}_{\PP^n_S}(a+r) \,\to \,
%   {\pi_S}_* \oplus^q {\mathcal O}_{\PP^n_S}(b+r) \,\to \,
%   {\pi_S}_* \F(r)                       \,\to \,
%   R^1{\pi_S}_* \G(r)$$
% and [...]
% There exists an integer $r_0$ such that
% $R^1{\pi_S}_*\G(r) =0$ and
% $R^1{\pi_T}_*\H(r)=0$ for all $r\ge r_0$. [...]
% The vertical maps are base-change homomorphisms, the first two
% of which are isomorphisms for all $r$.
% Therefore by the five lemma,
% $\phi^*{\pi_S}_* \F(r) \to {\pi_T}_* \F_T(r)$ is an isomorphism
% for all $r\ge r_0$."
\begin{proof}
  \uses{def:projective-space, thm:serre-vanishing, lem:five-lemma}
  Base change holds for open immersions, so using a finite affine open cover
  $\{U_i\}$ of $S$ and a finite affine open cover $\{V_{i,j}\}$ of each
  $\phi^{-1}(U_i)$ (possible by the noetherian hypothesis), we reduce to the case
  where $S = \spec A$ and $T = \spec B$ are affine.

  For every integer $i$ the base-change map
  $\phi^*{\pi_S}_*\Oo_{\PP^n_S}(i)\to {\pi_T}_*\Oo_{\PP^n_T}(i)$ is an
  isomorphism, and pull-back of a homomorphism $\Oo(a)\to\Oo(b)$ is compatible
  with these isomorphisms. Since $S$ is noetherian affine, there is a finite
  presentation by twists
  \[
    \textstyle\bigoplus^p \Oo_{\PP^n_S}(a)\xrightarrow{u}
    \bigoplus^q \Oo_{\PP^n_S}(b)\xrightarrow{v}\Ff\to 0,
  \]
  whose pull-back to $\PP^n_T$ is again exact. Put $\Gg=\ker(v)$ and
  $\Hh=\ker(v_T)$. Twisting and applying $\pi_*$ to the two short exact sequences
  yields, for every $r$, exact sequences ending in $R^1{\pi_S}_*\Gg(r)$ and
  $R^1{\pi_T}_*\Hh(r)$ respectively. By Serre vanishing
  (\cref{thm:serre-vanishing}) there is an $r_0$ with $R^1{\pi_S}_*\Gg(r)=0$ and
  $R^1{\pi_T}_*\Hh(r)=0$ for all $r\ge r_0$; hence both $\pi_*$-sequences become
  right-exact presentations of ${\pi_S}_*\Ff(r)$ and ${\pi_T}_*\Ff_T(r)$.
  Pulling the $S$-presentation back under $\phi$ (using right-exactness of
  $\ot$) gives a commutative diagram with exact rows whose first two vertical
  maps are the base-change isomorphisms above. By the five lemma
  (\cref{lem:five-lemma}) the third vertical map
  $\phi^*{\pi_S}_*\Ff(r)\to{\pi_T}_*\Ff_T(r)$ is an isomorphism for all
  $r\ge r_0$.
\end{proof}

\begin{remark}
  An alternative elementary proof (Mumford) identifies the graded
  $\Oo_S$-module $M=\bigoplus_{m}{\pi_S}_*\Ff(m)$ with $\Ff = M^\sim$, observes
  $\Ff_T=(\phi^*M)^\sim$ and $N=\bigoplus_m{\pi_T}_*\Ff_T(m)$ with
  $\Ff_T=N^\sim$, and deduces that the natural maps $(\phi^*M)_m\to N_m$ are
  isomorphisms for all $m\gg 0$.
\end{remark}

\begin{lemma}[Five lemma, right-exact form]
  \label{lem:five-lemma}
  \lean{LinearMap.bijective_of_surjective_of_bijective_of_right_exact}
  \mathlibok
  \textit{Provided by Mathlib.}
  Given a commutative diagram of modules with right-exact rows
  $M_1\to M_2\to M_3\to 0$ and $N_1\to N_2\to N_3\to 0$ and vertical maps
  $i_1,i_2,i_3$, if $i_1$ is surjective and $i_2$ is bijective then $i_3$ is
  bijective. (Used to conclude that the third base-change map is an isomorphism
  once the first two are.)
\end{lemma}

\section{Flatness of \texorpdfstring{$\Ff$}{F} from local freeness of \texorpdfstring{$\pi_*\Ff(r)$}{pi\_* F(r)}}

\begin{definition}\leanok
[Direct image commutes with base change]
  \label{def:direct-image-commutes-base-change}
  \lean{MR2223407ConstructionHilbertQuot.DirectImageCommutesBaseChange}
  A project-bespoke proxy predicate used to express ``$\Ff$ is flat over the
  base'' in the absence of a Mathlib sheaf-flatness notion: for $f\colon X\to S$
  and $\Ff$ on $X$, the proposition that formation of the underived direct image
  $\pi_*\Ff$ commutes with arbitrary base change $T\to S$. It is the operative
  consequence of flatness invoked downstream.
\end{definition}

\begin{lemma}\leanok
[Flatness from local freeness of direct images]
  \label{lem:flatness-from-local-freeness}
  \lean{MR2223407ConstructionHilbertQuot.flatness_of_locally_free_direct_image}
  \uses{def:projective-space, thm:coherent-higher-direct-image, def:direct-image-commutes-base-change}
  % SOURCE: Nitsure, §3, lines 1403--1408 (read from references/MR2223407-construction-hilbert-quot.tex)
  % SOURCE QUOTE: "Let $S$ be a noetherian scheme and let $\F$ be a coherent sheaf on
  % $\PP^n_S$. Suppose that there exists some integer $N$ such that
  % for all $r \ge N$ the direct image $\pi_*\F(r)$ is locally free.
  % Then $\F$ is flat over $S$."
  \textit{Source: Nitsure §3 (Flatness of $\Ff$ from Local Freeness of $\pi_*\Ff(r)$).}
  Let $S$ be a noetherian scheme and let $\Ff$ be a coherent sheaf on $\PP^n_S$.
  Suppose there exists an integer $N$ such that for all $r\ge N$ the direct image
  $\pi_*\Ff(r)$ is locally free. Then $\Ff$ is flat over $S$.
\end{lemma}

% SOURCE QUOTE PROOF: "Consider the graded module $M = \oplus_{r\ge N} M_r$ over
% ${\mathcal O}_S$, where $M_r =  \pi_*\F(r)$. The sheaf $\F$ is isomorphic
% to the sheaf $M^{\sim}$ on
% $\PP^n_S = {\bf Proj}_S\,{\mathcal O}_S[x_0,\ldots,x_n]$ made from the
% graded sheaf $M$ of ${\mathcal O}_S$-modules.
% As each $M_r$ is flat over ${\mathcal O}_S$, so is $M$. Therefore
% for any $x_i$ the localisation $M_{x_i}$ is flat over ${\mathcal O}_S$.
% There is a grading on $M_{x_i}$, indexed by $\Z$, defined
% by putting $\deg(v_p/x_i^q) = p-q$ for $v_p \in M_p$ (this is well-defined).
% Hence the component $(M_{x_i})_0$ of degree zero,
% being a direct summand of $M_{x_i}$, is again
% flat over ${\mathcal O}_S$. But by definition of $M^{\sim}$, this is just
% $\Gamma(U_i,\F)$, where $U_i = {\bf Spec}_S\,{\mathcal O}_S[x_0/x_i,\ldots,x_n/x_i]
% \subset \PP^n_S$. As the $U_i$ form an open cover of $\PP^n_S$,
% it follows that $\F$ is flat over ${\mathcal O}_S$."
\begin{proof}
  \uses{def:projective-space}
  Consider the graded $\Oo_S$-module $M=\bigoplus_{r\ge N}M_r$ with
  $M_r=\pi_*\Ff(r)$. Then $\Ff\cong M^\sim$ on
  $\PP^n_S=\Proj_S\Oo_S[x_0,\dots,x_n]$. Each $M_r$ is flat over $\Oo_S$ (locally
  free), hence so is $M$, and therefore so is each localization $M_{x_i}$. Grade
  $M_{x_i}$ over $\ZZ$ by $\deg(v_p/x_i^q)=p-q$ for $v_p\in M_p$; the degree-zero
  component $(M_{x_i})_0$ is a direct summand of $M_{x_i}$, hence flat over
  $\Oo_S$. By definition of $M^\sim$ this summand is $\Gamma(U_i,\Ff)$ on the
  standard affine $U_i=\Specc_S\Oo_S[x_0/x_i,\dots,x_n/x_i]\subset\PP^n_S$. As
  the $U_i$ cover $\PP^n_S$, the sheaf $\Ff$ is flat over $\Oo_S$.
\end{proof}

\begin{remark}
  The converse holds: if $\Ff$ is flat over $S$ then $\pi_*\Ff(r)$ is locally
  free for all sufficiently large $r$ (this is the statement quantified in
  \cref{thm:base-change-flat}).
\end{remark}

\section{The Grothendieck complex}

The following fundamental result of Grothendieck is the foundation for all
statements about direct images and base change of flat families; the complex
$K^\bullet$ occurring in it is called the \emph{Grothendieck complex}. It is the
local, explicit form of the base-change anchor \cref{thm:cohomology-base-change}
of Chapter~\ref{chap:basic}.

\begin{theorem}\leanok
[Grothendieck complex]
  \label{thm:grothendieck-complex}
  \lean{MR2223407ConstructionHilbertQuot.grothendieck_complex}
  \uses{thm:cohomology-base-change}
  % SOURCE: Nitsure, §3, lines 1459--1471 (read from references/MR2223407-construction-hilbert-quot.tex)
  % SOURCE QUOTE: "Let $\pi: X\to S$ be a proper morphism of noetherian
  % schemes where $S=\spec A$ is affine,
  % and let $\F$ be a coherent ${\mathcal O}_X$-module
  % which is flat over ${\mathcal O}_S$. Then
  % there exists a finite complex
  % $$0\to K^0 \to K^1 \to \ldots \to K^n\to 0$$
  % of finitely generated projective $A$-modules, together with
  % a functorial $A$-linear isomorphism
  % $$H^p(X, \F \otimes_A M)
  % \stackrel{\sim}{\to} H^p(K^{\cdot}\otimes_A M)$$
  % on the category of all $A$-modules $M$."
  \textit{Source: Nitsure §3 (Grothendieck Complex); EGA III 7.7.}
  Let $\pi\colon X\to S$ be a proper morphism of noetherian schemes with
  $S=\spec A$ affine, and let $\Ff$ be a coherent $\Oo_X$-module flat over
  $\Oo_S$. Then there exists a finite complex
  \[
    0\to K^0\to K^1\to\cdots\to K^n\to 0
  \]
  of finitely generated projective $A$-modules, together with a functorial
  $A$-linear isomorphism
  \[
    H^p(X,\Ff\ot_A M)\;\xrightarrow{\ \sim\ }\;H^p(K^\bullet\ot_A M)
  \]
  on the category of all $A$-modules $M$.
\end{theorem}

\begin{proof}
  \uses{thm:cohomology-base-change}
  \leanok
  This is the affine-base, explicit form of the cohomology-and-base-change anchor
  \cref{thm:cohomology-base-change}. It is a classical theorem of Grothendieck;
  see EGA III 7.7 (and Mumford, \emph{Abelian Varieties}, §5, for the standard
  exposition). The complex $K^\bullet$ is obtained from the \v{C}ech complex of a
  finite affine cover of $X$ by replacing it with a quasi-isomorphic complex of
  finite projective $A$-modules, using flatness of $\Ff$ to ensure base-change
  compatibility. Treated as a foundational input; not reproved here.
\end{proof}

As a consequence we have the following representability statement (EGA III
7.7.6), which produces the coherent sheaf used to represent $\Hom$-functors in
the Quot construction.

\begin{theorem}\leanok
[Cokernel of the dual {\rm (EGA III 7.7.6)}]
  \label{thm:cokernel-of-dual}
  \lean{MR2223407ConstructionHilbertQuot.cokernel_of_dual}
  \uses{thm:grothendieck-complex}
  % SOURCE: Nitsure, §3, lines 1481--1490 (read from references/MR2223407-construction-hilbert-quot.tex)
  % SOURCE QUOTE: "Let $S$ be a noetherian scheme and $\pi : X \to S$ a proper morphism.
  % Let $\F$ be a coherent sheaf on $X$ which is flat over $S$.
  % Then there exists a coherent sheaf $\Qq$ on $S$ together with
  % a functorial ${\mathcal O}_S$-linear isomorphism
  % $$\theta_{\G} : \pi_*(\F \otimes_{{\mathcal O}_X} \pi^*\G) \to \hom_{{\mathcal O}_S}(\Qq,\G)$$
  % on the category of all quasi-coherent sheaves $\G$ on $S$.
  % By its universal property,
  % the pair $(\Qq, \theta)$ is unique up to a unique isomorphism."
  \textit{Source: Nitsure §3 (EGA III 7.7.6).}
  Let $S$ be a noetherian scheme and $\pi\colon X\to S$ a proper morphism. Let
  $\Ff$ be a coherent sheaf on $X$ which is flat over $S$. Then there exists a
  coherent sheaf $\mathcal{Q}$ on $S$ together with a functorial
  $\Oo_S$-linear isomorphism
  \[
    \theta_{\Gg}\colon \pi_*\bigl(\Ff\ot_{\Oo_X}\pi^*\Gg\bigr)
      \;\longrightarrow\;\mathcal{H}om_{\Oo_S}(\mathcal{Q},\Gg)
  \]
  on the category of all quasi-coherent sheaves $\Gg$ on $S$. By its universal
  property the pair $(\mathcal{Q},\theta)$ is unique up to unique isomorphism.
  When $\mathcal{Q}$ is used to form the affine $S$-scheme
  $\spec\Sym_{\Oo_S}\mathcal{Q}$ (a \emph{linear scheme}), this representability
  is what makes the relevant $\Hom$-functors representable downstream.
\end{theorem}

% SOURCE QUOTE PROOF: "If $S =\spec A$, then
% we can take $\Qq$ to be the coherent sheaf associated to the $A$-module
% $Q$ which is the cokernel of the
% transpose $\pa^{\vee} : (K^1)^{\vee} \to (K^0)^{\vee}$ where
% $\pa: K^0 \to K^1$ is the differential of any chosen Grothendieck
% complex of $A$-modules
% $0\to K^0 \to K^1 \to \ldots \to K^n\to 0$ for the
% sheaf $\F$, whose existence is given by Theorem \ref{Grothendieck complex}.
% For any $A$-module $M$,
% the right-exact sequence
% $(K^1)^{\vee} \to (K^0)^{\vee} \to Q\to 0$ with $M$
% gives on applying $Hom_A(-,M)$ a left-exact sequence
% $$0 \to Hom_A(Q, M) \to K^0\otimes_AM \to K^1\otimes_AM$$
% Therefore by Theorem \ref{Grothendieck complex}, we have
% an isomorphism
% $$\theta^A_M : H^0(X_A, \F_A \otimes_A M)  \to  Hom_A(Q, M)$$
% Thus, the pair $(Q, \theta^A)$
% satisfies the theorem when $S = \spec A$. More generally, we can cover
% $S$ by affine open subschemes. Then on their overlaps,
% the resulting pairs $(\Qq,\theta)$ glue together by their uniqueness."
\begin{proof}
  \uses{thm:grothendieck-complex}
  When $S=\spec A$, choose a Grothendieck complex
  $0\to K^0\xrightarrow{\partial}K^1\to\cdots\to K^n\to 0$ for $\Ff$
  (\cref{thm:grothendieck-complex}), and let $Q=\coker\bigl(\partial^\vee\colon
  (K^1)^\vee\to(K^0)^\vee\bigr)$ with $\mathcal{Q}=Q^\sim$. For any $A$-module
  $M$, applying $\Hom_A(-,M)$ to the right-exact sequence
  $(K^1)^\vee\to(K^0)^\vee\to Q\to 0$ gives a left-exact sequence
  $0\to\Hom_A(Q,M)\to K^0\ot_A M\to K^1\ot_A M$, whose kernel is by
  \cref{thm:grothendieck-complex} identified with $H^0(X_A,\Ff_A\ot_A M)$. This
  furnishes the isomorphism $\theta^A_M$. For general $S$, the locally
  constructed pairs $(\mathcal{Q},\theta)$ glue over an affine cover by their
  uniqueness.
\end{proof}

\section{Semicontinuity and base change for flat sheaves}

The following is Grothendieck's main theorem on base change for flat families,
a consequence of the Grothendieck complex (\cref{thm:grothendieck-complex}). For
clarity we record its semicontinuity content and its base-change content as two
statements; both are part of one theorem in EGA III 7.7.8--7.7.9.

\begin{theorem}\leanok
[Upper semicontinuity of fibre cohomology {\rm (EGA III 7.7.8--7.7.9)}]
  \label{thm:semicontinuity}
  \lean{MR2223407ConstructionHilbertQuot.semicontinuity_fibre_cohomology}
  \uses{thm:grothendieck-complex}
  % SOURCE: Nitsure, §3, lines 1618--1635 (read from references/MR2223407-construction-hilbert-quot.tex)
  % SOURCE QUOTE: "Let $\pi: X\to S$ be a proper morphism of noetherian
  % schemes, and let $\F$ be a coherent ${\mathcal O}_X$-module
  % which is flat over ${\mathcal O}_S$.
  % Then the following statements hold:
  % (1) For any integer $i$ the function $s\mapsto \dim_{\kappa(s)}H^i(X_s,\F_s)$
  % is upper semi-continuous on $S$,
  % (2) The function $s\mapsto \sum_i (-1)^i \dim_{\kappa(s)}H^i(X_s,\F_s)$
  % is locally constant on $S$.
  % (3) If for some integer $i$, there is some integer $d\ge 0$ such that
  % for all $s\in S$ we have $\dim_{\kappa(s)}H^i(X_s,\F_s) = d$,
  % then $R^i\pi_*\F$ is locally free of rank $d$,
  % and $(R^{i-1}\pi_*\F)_s \to H^{i-1}(X_s,\F_s)$ is an isomorphism
  % for all $s\in S$."
  \textit{Source: Nitsure §3 (Base-change for flat sheaves, parts (1)--(3); EGA III 7.7.8--7.7.9).}
  Let $\pi\colon X\to S$ be a proper morphism of noetherian schemes and $\Ff$ a
  coherent $\Oo_X$-module flat over $\Oo_S$. Then:
  \begin{enumerate}
    \item for every integer $i$ the function
      $s\mapsto\dim_{\kappa(s)}H^i(X_s,\Ff_s)$ is upper semicontinuous on $S$;
    \item the function $s\mapsto\sum_i(-1)^i\dim_{\kappa(s)}H^i(X_s,\Ff_s)$ is
      locally constant on $S$;
    \item if for some $i$ there is $d\ge 0$ with
      $\dim_{\kappa(s)}H^i(X_s,\Ff_s)=d$ for all $s\in S$, then $R^i\pi_*\Ff$ is
      locally free of rank $d$ and $(R^{i-1}\pi_*\Ff)_s\to H^{i-1}(X_s,\Ff_s)$ is
      an isomorphism for all $s\in S$.
  \end{enumerate}
\end{theorem}

% SOURCE QUOTE PROOF: "See for example Hartshorne [H] Chapter III, Section 12 for a proof.
% It is possible to replace the use of the formal function
% theorem in [H] (or the original argument in [EGA] based on
% completions) in proving the statement (4) above,
% with an elementary argument based on applying
% Nakayama lemma to the Grothendieck complex."
\begin{proof}
  \uses{thm:grothendieck-complex}
  All assertions are local on $S$, so we may assume $S=\spec A$ and use a
  Grothendieck complex $K^\bullet$ of finite projective $A$-modules computing the
  cohomology of $\Ff$ universally (\cref{thm:grothendieck-complex}). Then
  $\dim_{\kappa(s)}H^i(X_s,\Ff_s)=\dim_{\kappa(s)}H^i(K^\bullet\ot\kappa(s))$,
  which jumps up on closed conditions (ranks of the differentials drop), giving
  (1); the alternating sum equals $\sum(-1)^p\rank K^p$, giving (2); and (3) is
  the rigidity of the complex when a fibre dimension is constant. See Hartshorne
  III.12 (or EGA III §7) for the standard argument.
\end{proof}

\begin{theorem}\leanok
[Base change for flat sheaves {\rm (EGA III 7.7.8--7.7.9)}]
  \label{thm:base-change-flat}
  \lean{MR2223407ConstructionHilbertQuot.base_change_for_flat_sheaves}
  \uses{thm:grothendieck-complex, lem:flatness-from-local-freeness}
  % SOURCE: Nitsure, §3, lines 1636--1657 (read from references/MR2223407-construction-hilbert-quot.tex)
  % SOURCE QUOTE: "(4) If for some integer $i$ and point $s\in S$ the map
  % $(R^i\pi_*\F)_s \to H^i(X_s,\F_s)$ is surjective,
  % then there exists
  % an open subscheme $U\subset S$ containing $s$ such that
  % for any quasi-coherent ${\mathcal O}_U$-module $\G$ the natural homomorphism
  % $$(R^i {\pi_U}_*\F_{X_U})\otimes_{{\mathcal O}_U}\G\to
  % R^i{\pi_U}_*(\F_{X_U}\otimes_{{\mathcal O}_{X_U}}{\pi_U}^*\G)$$
  % is an isomorphism, where $X_U = \pi^{-1}(U)$ and
  % $\pi_U: X_U \to U$ is induced by $\pi$. In particular,
  % $(R^i\pi_*\F)_{s'} \to H^i(X_{s'},\F_{s'})$ is an isomorphism
  % for all $s'$ in $U$.
  % (5) If for some integer $i$ and point $s\in S$ the map
  % $(R^i\pi_*\F)_s \to H^i(X_s,\F_s)$ is surjective,
  % then the following conditions (a) and (b) are equivalent:
  % ~~~~(a) The map $(R^{i-1}\pi_*\F)_s \to H^{i-1}(X_s,\F_s)$ is surjective.
  % ~~~~(b) The sheaf $R^i\pi_*\F$ is locally free in a
  %         neighbourhood of $s$ in $S$."
  \textit{Source: Nitsure §3 (Base-change for flat sheaves, parts (4)--(5); EGA III 7.7.8--7.7.9).}
  With $\pi\colon X\to S$ proper between noetherian schemes and $\Ff$ coherent
  and flat over $S$:
  \begin{enumerate}
    \item[(4)] if for some $i$ and $s\in S$ the map
      $(R^i\pi_*\Ff)_s\to H^i(X_s,\Ff_s)$ is surjective, then there is an open
      $U\ni s$ such that for every quasi-coherent $\Oo_U$-module $\Gg$ the
      natural map
      $(R^i{\pi_U}_*\Ff_{X_U})\ot_{\Oo_U}\Gg\to
       R^i{\pi_U}_*(\Ff_{X_U}\ot_{\Oo_{X_U}}\pi_U^*\Gg)$
      is an isomorphism; in particular $(R^i\pi_*\Ff)_{s'}\to H^i(X_{s'},\Ff_{s'})$
      is an isomorphism for all $s'\in U$;
    \item[(5)] under the same surjectivity hypothesis, $R^i\pi_*\Ff$ is locally
      free near $s$ iff $(R^{i-1}\pi_*\Ff)_s\to H^{i-1}(X_s,\Ff_s)$ is surjective.
  \end{enumerate}
  In particular, applied with $i=0$ to the twists $\Ff(r)$ for $X=\PP^n_S$ and
  $r\gg 0$ (where the higher cohomology of the flat $\Ff$ vanishes fibrewise),
  $\pi_*\Ff(r)$ is locally free and commutes with arbitrary base change; this is
  the converse to \cref{lem:flatness-from-local-freeness}.
\end{theorem}

\begin{proof}
  \uses{thm:grothendieck-complex, lem:flatness-from-local-freeness}
  Work locally with a Grothendieck complex $K^\bullet$ as in
  \cref{thm:grothendieck-complex}. Surjectivity of
  $(R^i\pi_*\Ff)_s\to H^i(X_s,\Ff_s)$ is, via the complex, the statement that the
  $i$-th cohomology of $K^\bullet$ is computed compatibly with $\ot\kappa(s)$;
  Nakayama's lemma applied to $K^\bullet$ then propagates this to a neighbourhood
  and yields the universal base-change isomorphism in (4), replacing the formal
  function theorem used in Hartshorne III.12. Part (5) is the local-freeness
  criterion for $R^i\pi_*\Ff$ in terms of exactness of the dual complex at the
  relevant spot. For $i=0$ and $r\gg 0$, fibrewise Serre vanishing makes the
  surjectivity hypothesis automatic, so $\pi_*\Ff(r)$ is locally free and
  base-change compatible; together with \cref{lem:flatness-from-local-freeness}
  this is the announced converse. See Hartshorne III.12 / EGA III §7.
\end{proof}
